\documentclass[12pt]{article}
\usepackage[utf8]{inputenc}
\usepackage[T1]{fontenc}
\usepackage{hyperref}
\title{Tsuvadi Language Resources for K-12 Curriculum Development}
\author{Compiled by Claude Sonnet 5 (claude-sonnet-5) for Andr\'as Kornai}
\date{}
\begin{document}
\maketitle

\section{Introduction}
Tsuvadi (ISO 639-3 tvd; also Vadi, Abadi, Avadi, Evadi, Kamberi) is a Kainji language of Nigeria spoken by the Kambari people, mainly in Kebbi and Niger states. Omniglot reports about 130,000 speakers and says the language is used by all ages in daily life, except in mosques. Search results turned up little material that is usable in schools. This report says so plainly and does not pad.

\section{Primary Educational Resources}
No primers, textbooks or school programmes for Tsuvadi were found in the search results. Omniglot notes that few people can read and write the language.
\begin{itemize}
\item Omniglot overview of Vadi (TsuVa\d{d}i), with orthographic notes: \url{https://www.omniglot.com/writing/vadi.htm}
\end{itemize}

\section{Linguistic and Reference Materials}
\begin{itemize}
\item Wikipedia article on the Vadi language. It cites Buttacavoli (2022), a phonological sketch of the Tsogudu dialect, and Blench (2007): \url{https://en.wikipedia.org/wiki/Vadi_language}
\item Blench 2007 phonology and wordlist source, listed in PHOIBLE: \url{https://phoible.org/sources/1653_tvd_blench}
\item Endangered Languages catalogue entry: \url{https://www.endangeredlanguages.com/bn/node/96036}
\end{itemize}

\section{Religious and Community Resources}
Joshua Project reports recent efforts to translate Scripture into the language. Omniglot says parts of the Bible have been translated.
\begin{itemize}
\item Joshua Project, Avadi/Tsivadi in Nigeria: \url{https://joshuaproject.net/people_groups/15595/NI}
\end{itemize}

\section{Digital and Audio Resources}
\begin{itemize}
\item Global Recordings Network, audio recordings in the Shambo variety: \url{https://globalrecordings.net/en/language/2079}
\end{itemize}

\section{Curriculum Development Context}
The language is vigorous in the community, which is a strong basis for mother-tongue instruction. However, no standardized orthography materials, readers or teacher resources were found. A curriculum would need to start from the published phonological work, the wordlist and community-led orthography development. The dialects (for example Tsogudu, Kakihum, Agadi) differ, so a variety must be chosen.
\end{document}